% =====================================================================
%  CAPÍTULO I / 2 Fundamentação Teórica
% =====================================================================
\setdiv{CAPÍTULO I}
\chapter{Fundamentação Teórica}

\section{Macro, Microeconomia e Inflação}

O cenário econômico global recente tem sido impactado por choques simultâneos:
conflitos geopolíticos, problemas gerais nas cadeias produtivas e pressões
inflacionárias de alto risco. Conforme o \citeonline{fmi2024}, esse cenário
evidencia um problema estrutural do sistema internacional: a liquidez global supera,
em alguns níveis, a produção real, criando distorções gerais nos preços e no
crescimento econômico.

No Brasil, a microeconomia sofre os efeitos desse ambiente. A taxa Selic em 14,25\%
\cite{bcbSelic} desestimula investimentos em produção real, como comprovam os estudos
do \citeonline{bcbJuros}, que relacionam juros altos à migração do capital para
aplicações financeiras, reduzindo posições no setor industrial.

A inflação, por fim, é um desafio em especial. O \citeonline{fmi2024} alertou que as
economias emergentes, assim como a brasileira, enfrentam o ``trade-off'' perverso:
políticas monetárias com grande enfoque restritivo controlam os preços, mas
sacrificam, no curto prazo, o crescimento do PIB.

Diante desse cenário geral, alternativas economicamente viáveis tornam-se uma
prioridade, especialmente as menos atreladas às variáveis macroeconômicas. Essa
contextualização servirá de pano de fundo para o estudo de estratégias de análise
econômica, que serão detalhadas no decorrer do projeto.

\section{Corretora}

De acordo com o BCB (Banco Central do Brasil), as corretoras e distribuidoras, na
atividade de intermediação, oferecem serviços como plataformas de investimento pela
internet (home broker), consultoria financeira, clubes de investimentos,
financiamento para compra de ações (conta margem) e administração e custódia de
títulos e valores mobiliários dos clientes. Na remuneração pelos serviços, essas
instituições podem cobrar comissões e taxas \cite{bcbCorretoras}.

Portanto, a substancial necessidade de um serviço mediador para representação em
ativos de capital aberto, ou seja, ativos que possuem participação em mercado
anunciadas para o acesso do grande público. Além disso, existem séries de medidas
formais que se demonstram necessárias para o processo burocrático da parte fiscal e
jurídica na transação de participações em empresas, medidas essas que são assumidas
por serviços de corretagem.

\section{Bolsa de Valores}

De acordo com o BCB (Banco Central do Brasil), as bolsas de valores são sociedades
anônimas ou associações civis, com o objetivo de manter local ou sistema adequado ao
encontro de seus membros e à realização entre eles de transações de compra e venda
de títulos e valores mobiliários, em mercado livre e aberto, especialmente
organizado e fiscalizado por seus membros e pela Comissão de Valores Mobiliários.
Possuem autonomia financeira, patrimonial e administrativa (Resolução CMN 2.690, de
2000) \cite{bcbBolsas}.

Diante da necessidade de capital exógeno, as empresas possuem o mecanismo de capital
aberto que possibilita a entrada de capital de fora da empresa por parte de
investidores; em contrapartida, a empresa cede porcentagens do negócio proporcionais
aos montantes investidos por parte dos acionistas. Concluindo, a prática é
extremamente relevante para fomentar e financiar o escalonamento da empresa, mesmo
que sejam cedidos percentuais do empreendimento, o montante geral escalonado é,
geralmente, maior que o montante de capital fechado.

\section{Ativos}

De acordo com a Warren Magazine, os ativos são todos aqueles que possuem expectativa
de superar o desempenho do mercado ou de um índice de referência. Utilizando-se de
análises detalhadas, pesquisas de mercado e previsões econômicas, os gestores se
alinham para decisões por posições em investimentos \cite{warren2025}.

Por fim, os gestores buscam por oportunidades de investimento que possam gerar
retornos acima da média do ativo e do mercado, alinhando-se ao princípio de alocação
eficiente de capital. Assim como destacado por \citeonline{elder2002}, cada decisão
de investimento requer uma análise minuciosa do trade-off, escolha com custo de
oportunidade, entre risco e retorno, uma vez que os recursos são limitados e devem
ser distribuídos de forma otimizada. Todavia, em cenários adversos, pode-se demonstrar
necessário o rebalanceamento da carteira, estrategicamente mitigando perdas de
capital, prática comum em mercados voláteis, assim como foi apontado por
\citeonline{ayres2025}.

\section{Tipos de Análise Técnica}

Uma forma de tentar prever o que vai acontecer com os preços dos ativos no mercado
financeiro é usar a análise técnica. Essa análise estuda o histórico dos preços para
encontrar padrões e tendências que possam indicar o que pode acontecer no futuro.
Essa ideia surgiu há muito tempo, com o comércio de arroz, e hoje é usada para
analisar ações, moedas e outros ativos. O importante é olhar para os preços, sem se
preocupar com notícias ou outras coisas que possam influenciar o mercado
\cite{fernandes2014}.

A análise técnica pode ser usada em quase todos os mercados e mistura ciência com a
forma como as pessoas pensam e agem nesses mercados. Para ter mais chances de
sucesso, é bom se especializar em um mercado e escolher as ferramentas certas para
analisar os dados. Existem dois jeitos principais de fazer análise técnica: a mais
antiga, chamada empírica ou clássica, e a mais moderna, chamada estatística
\cite{fernandes2014}.

\subsection{Análise Técnica Empírica ou Clássica}
A análise técnica empírica, ou clássica, é o método mais tradicional de analisar o
mercado financeiro. Nessa abordagem, os analistas usavam seus gráficos impressos,
observando cuidadosamente o movimento dos preços manualmente, com ferramentas mais
simples e rudimentares como réguas e lápis \cite{fernandes2014}.

O objetivo dela é identificar tendências de alta ou de baixa nos preços, assim como
encontrar áreas onde os preços geralmente encontram dificuldades para ultrapassar
(resistências) ou onde encontram apoio e param de cair (suportes). Além disso, esse
analista busca por padrões gráficos, que se repetem ao longo do tempo e que, segundo
essa abordagem, podem sinalizar movimentos futuros dos preços \cite{fernandes2014}.

Apesar de sua tradição, essa análise pode ser um tanto subjetiva. A interpretação dos
gráficos, a identificação dos padrões e a definição das áreas de suporte e resistência
dependem muito da experiência e da percepção de cada analista, o que pode levar a
conclusões diferentes sobre o mesmo ativo \cite{fernandes2014}.

\subsection{Análise Técnica Estatística ou Moderna}
A análise técnica estatística, ou moderna, representa uma abordagem mais científica
para analisar o mercado financeiro. Em vez de depender da interpretação visual e
subjetiva dos gráficos, os analistas utilizam modelos matemáticos e estatísticos para
analisar os preços dos ativos \cite{fernandes2014}.

Para isso, utilizam programas de computador que calculam uma variedade de indicadores
técnicos, como médias móveis (que suavizam os movimentos dos preços e ajudam a
identificar tendências), osciladores (como o Índice de Força Relativa -- RSI, que
indica se um ativo está ``sobrecomprado'' ou ``sobrevendido'') e muitos outros. Esses
indicadores fornecem informações objetivas sobre o comportamento do mercado,
auxiliando na tomada de decisões \cite{fernandes2014}.

Essa abordagem é considerada mais objetiva, pois minimiza a influência da
subjetividade na interpretação dos dados, utilizando métodos quantitativos e modelos
matemáticos para auxiliar na tomada de decisões \cite{fernandes2014}.

Neste trabalho, adotaremos a análise técnica estatística para criar uma ferramenta
que utilize modelos avançados para identificar oportunidades de negociação no
mercado. Essa abordagem se adequa à utilização de modelos complexos e algoritmos que
buscam aumentar a precisão e a eficiência na identificação de padrões e na avaliação
de riscos, permitindo a criação de estratégias de negociação mais sofisticadas e
adaptáveis às diferentes condições do mercado.

\section{MetaTrader 5}

O MetaTrader 5 (MT5) é plataforma de negociação multimercado avançada, projetada para
atender às crescentes demandas dos mercados financeiros globais. Oferecendo um
ambiente integrado para a negociação de Forex, ações, futuros e commodities, o MT5
proporciona uma experiência completa, que abrange desde a execução de ordens até a
análise técnica avançada e o desenvolvimento de sistemas automatizados
\cite{metaquotes2024}.

\begin{figure}[htb]
\centering
\caption{Interface MetaTrader}
\label{fig:interface-mt}
\includegraphics[width=0.85\textwidth]{fig01-interface-metatrader.png}
\fonte{\citeonline{metaquotes2024}.}
\end{figure}

A plataforma MetaTrader 5 (MT5), como se observa na imagem acima, oferece uma
interface personalizável que permite aos usuários organizarem múltiplas janelas de
acordo com suas preferências. A janela de Cotações exibe os preços em tempo real dos
instrumentos financeiros selecionados, e os Gráficos de Preços permitem visualizar a
evolução dos preços ao longo do tempo. O Navegador oferece acesso rápido a contas,
indicadores técnicos e scripts, facilitando a análise e o acompanhamento do mercado
\cite{metaquotes2024}.

\begin{figure}[htb]
\centering
\caption{Fluxograma do Funcionamento do EA}
\label{fig:fluxograma-ea}
\includegraphics[width=0.55\textwidth]{fig02-fluxograma-ea.png}
\fonte{\citeonline{almeida2016}.}
\end{figure}

% TODO: citação (FLÁVIO, 2023) removida das duas frases abaixo — não há
% entrada correspondente em referencias.bib e a fonte não foi localizada.
% Confirmar a referência real e reinserir com \citeonline, ou manter sem
% citação caso a afirmação seja considerada conhecimento geral da área.
Um dos recursos mais poderosos do MT5 são os \textit{Expert Advisor}s (EAs), também conhecidos
como robôs de negociação. EAs são programas que automatizam o processo
de negociação, permitindo que os usuários executem estratégias complexas sem a
necessidade de intervenção manual. Esses robôs abrem, modificam e
fecham ordens automaticamente, com base em algoritmos predefinidos, e podem ser
criados tanto manualmente, para maior customização, quanto automaticamente, por meio
do MQL5 \cite{almeida2016}.

A análise técnica é um componente fundamental do MT5, com um conjunto abrangente de
ferramentas para a identificação de padrões e sinais de mercado. A plataforma oferece
uma variedade de indicadores técnicos integrados, abrangendo todas as principais
categorias, bem como objetos gráficos e ferramentas de desenho para uma análise
visual detalhada dos gráficos. A capacidade de analisar dados em múltiplos períodos,
desde M1 (1 minuto) até MN1 (1 mês), permite capturar padrões que podem não ser
evidentes em uma única escala temporal, resultando em algoritmos de previsão mais
robustos \cite{metaquotes2024}.

A negociação algorítmica é um dos pilares do MT5, devido à linguagem MQL5, que
permite criar EAs para a análise automatizada de dados de mercado, a geração de sinais
de negociação e a execução de ordens. O MetaEditor oferece um ambiente de
desenvolvimento completo, com recursos como destaque de sintaxe, compilador integrado
e depurador. O Testador de Estratégias possibilita a validação e otimização de
algoritmos utilizando dados históricos, com opções para testar em diferentes períodos
para simulação precisa e analisar estatísticas de desempenho detalhadas
\cite{metaquotes2024}.

A integração com Python, através de uma API dedicada, amplia as capacidades do MT5,
permitindo a utilização de bibliotecas de aprendizado de máquina e a criação de
modelos preditivos sofisticados. Essa integração possibilita a combinação das
funcionalidades nativas da plataforma com o poder analítico do ecossistema Python,
resultando em algoritmos de previsão mais precisos e adaptáveis. Adicionalmente, o
MT5 oferece suporte nativo para redes neurais em MQL5 e computação paralela com
OpenCL, fornecendo recursos avançados para o desenvolvimento de algoritmos
computacionalmente intensivos \cite{metaquotes2024}.

Em suma, o MetaTrader 5 se consolida como uma plataforma versátil e poderosa para o
desenvolvimento e a implementação de algoritmos de previsão de mercado. Sua
combinação de recursos nativos, integração com tecnologias externas e ecossistema
maduro o tornam uma escolha estratégica para pesquisadores, desenvolvedores e traders
que buscam criar sistemas preditivos sofisticados e eficientes para os mercados
financeiros.

\section{Indicadores}

Indicadores Técnicos são ferramentas fundamentais na análise estatística de mercados
financeiros, pois traduzem fórmulas matemáticas e gráficos de forma que auxiliam os
traders na tomada de decisões. Utilizando dados históricos de preço, volume, ou a
combinação de ambos, os indicadores buscam revelar padrões e tendências que, de outra
forma, seriam difíceis de discernir. Seu propósito principal é identificar potenciais
pontos de entrada e saída em negociações, sinalizar a força ou fraqueza de uma
tendência estabelecida, e demarcar zonas de suporte e resistência onde o preço tende
a encontrar dificuldades em romper \cite{fernandes2014}.

Embora baseados em cálculos estatísticos, os sinais gerados por indicadores técnicos
não devem ser interpretados como previsões infalíveis. Pelo contrário, eles funcionam
como alertas ou confirmações, que exigem uma análise cuidadosa e contextualizada por
parte do trader. A eficácia de um indicador pode variar significativamente entre
diferentes ativos e ao longo do tempo, exigindo um processo contínuo de adaptação e
ajuste de parâmetros \cite{fernandes2014}.

Em seguimento, alguns indicadores específicos serão explorados neste trabalho,
visando a identificação de tendências e a confirmação de sinais de negociação. A
escolha e a aplicação desses indicadores serão detalhadas nas seções subsequentes.

\subsection{Índice de Força Relativa}
De acordo com \citeonline{anderson1993}, o RSI é um dos mais conhecidos indicadores
de análise técnica e amplamente utilizado, não somente no mercado financeiro, mas
para medir a consolidação de uma tendência em diversas áreas. O índice é base para
diversos outros índices que necessitam de um norte para onde a tendência do ativo
aponta, sendo assim, uma ferramenta imprescindível para fomentar uma decisão
estratégica de entrada e saída de um ativo com maior segurança.

Diante do exposto, o Índice de Força Relativa tem como objetivo estudar posições
estratégicas no preço dos ativos para aproveitar-se de uma brecha de compra ou venda.

Sendo assim, as classificações de ativos por RSI são basicamente sobrevalorizados ou
supervalorizados; quando supervalorizado a melhor posição no ativo é de venda pois o
ativo vem esgotando sua força compradora devido ao preço inflado, em contrapartida, o
sobrevalorizado indica uma posição forte de compra pois, devido ao preço baixo, não
será recompensador para o investidor vender o ativo, criando-se uma resistência que
não propicia a vendas e sim compras.

\begin{figure}[htb]
\centering
\caption{Tendência de Alta com RSI}
\label{fig:rsi}
\includegraphics[width=0.8\textwidth]{fig03-tendencia-alta-rsi.png}
\fonte{\citeonline{fernandes2014}.}
\end{figure}

As margens de operação do RSI para compra e venda são frequentemente revisitadas; a
estrutura clássica diz que, abaixo de 30 pontos no índice, o ativo tende a subir pois
esgotou sua força vendedora, ao mesmo passo que, acima de 70 pontos no índice, o ativo
tende a cair pois esgotou sua força compradora. Levando em consideração todo o
contexto, o RSI não representa uma determinante que afirma o momento de compra ou
venda, o índice tem como objetivo ser uma ferramenta a mais para o consenso de
entradas e saídas em investimentos.

\subsection{Médias Móveis}
As Médias Móveis (MM) são um dos indicadores técnicos mais amplamente utilizados na
análise estatística de mercados financeiros, oferecendo uma maneira eficaz de suavizar
as flutuações de preço e identificar tendências predominantes. Ao calcular o preço
médio de um ativo durante um período específico, as Médias Móveis reduzem o ruído
gerado por movimentos de preço aleatórios, permitindo que os analistas visualizem a
direção geral do mercado com maior clareza, fornecendo pontos de referência válidos ao
longo do tempo, possibilitando a comparação entre níveis atuais e passados
\cite{fernandes2014}.

Existem diferentes tipos de Médias Móveis, sendo as mais comuns a Média Móvel Simples
(MMS) e a Média Móvel Exponencial (MME). A MMS atribui o mesmo peso a todos os preços
dentro do período, enquanto a MME pondera os preços mais recentes de forma mais
significativa, tornando-a mais sensível às mudanças de preço. A escolha entre MMS e
MME depende do objetivo da análise e da sensibilidade desejada. Períodos mais curtos
resultam em Médias Móveis mais sensíveis, que reagem rapidamente às mudanças de preço,
mas também geram mais sinais falsos. Períodos mais longos, por outro lado, resultam em
Médias Móveis mais suaves, que filtram o ruído, mas podem atrasar a identificação de
novas tendências \cite{fernandes2014}.

\begin{figure}[htb]
\centering
\caption{Cruzamento de médias móveis}
\label{fig:cruzamento-mm}
\includegraphics[width=0.95\textwidth]{fig04-cruzamento-medias.png}
\fonte{Elaboração própria com dados e interface do MetaTrader 5.}
\end{figure}

Uma aplicação comum das Médias Móveis, como é possível observar na imagem acima, é a
identificação de cruzamentos, interpretados como sinais de compra ou venda.
\citeonline{almeida2016}, em sua pesquisa, utilizou essa técnica com duas médias, uma
``rápida'' e outra ``lenta'', gerando sinais de negociação com base em seus
cruzamentos. Para um controle de risco mais refinado, implementou também StopLoss e
TakeProfit. Entretanto, \citeonline{almeida2016} alerta para as limitações de se usar
apenas essa técnica, especialmente em mercados com alta volatilidade, reforçando a
necessidade de ferramentas que confirmem o sinal para evitar falsos positivos.

\subsection{Médias Móveis Exponenciais}
As Médias Móveis Exponenciais (MME) constituem uma ferramenta amplamente conhecida na
análise técnica geral, destacando-se no mesmo sentido das médias móveis (MM), sendo
principalmente utilizada para suavizar alterações bruscas de curto prazo nos preços
dos ativos; além disso, atribui pesos mais acentuados para dados recentes. A diferença
entre os indicadores MM e MME é fundamentalmente velocidade: devido à fórmula adaptada
da MME, os resultados são mais velozes nas comparações utilizadas, contudo, o
processamento requisitado é maior. Por esse motivo, a MME é significativamente
utilizada em mercados de volatilidade pois é de suma importância identificar padrões
de reversões ou de continuidade num espaço de tempo consideravelmente curto para
tomadas de decisões.

Como já afirmado, a relevância da MME é ampla na área de análise técnica.
\citeonline{elder2002}, por exemplo, destaca o seu valor para análise de um mercado
atual, característica essa que é contrária em outros indicadores onde suas fórmulas não
consideram as datas dos dados extraídos. Segundo o autor, ``a média móvel exponencial
permite que o trader veja o mercado como ele é agora, e não como era semanas atrás''
\cite[p.~145]{elder2002}. Seguindo sua linha de pensamento, o cruzamento de médias
possui a capacidade de demonstrar sinais de entrada ou saída de forma eficaz.

Apesar de tantas vantagens, a MME também possui seus pontos de consideração. A
principal característica da MME é resultante da sua velocidade em extrair informações
que, em mercados de lateralização, gera sinais que podem ser interpretados como
falsos, levando a entradas e saídas em momentos desinteressantes. Uma forma de
diminuir esta influência é a complementação com outros indicadores, como por exemplo,
Índice de Força Relativa (IFR) ou Bandas de Bollinger; essa complementação cria dados
cruzados que confirmam uma incerteza ou uma compra, facilitando decisões.

\subsection{Bandas de Bollinger}
As Bandas de Bollinger (BB), um indicador técnico de volatilidade desenvolvido por
John Bollinger, oferecem uma perspectiva sobre a dinâmica dos preços em um mercado. As
BB são construídas a partir de uma Média Móvel Aritmética (MM20, por exemplo) aplicada
aos preços, traçando duas curvas que representam os limites superior e inferior da
volatilidade esperada para o ativo. Essas bandas são calculadas utilizando o desvio
padrão da média móvel, de modo que quanto maior a volatilidade do ativo, maior será a
distância entre as bandas \cite{fernandes2014}.

Ao contrário de uma interpretação simplista que considera a banda superior como uma
resistência e a banda inferior como um suporte, a teoria de Bollinger enfatiza a
importância de analisar o comportamento dos preços em relação aos seus limites de
volatilidade esperados. Se, em um movimento de alta, o preço atinge a banda superior,
isso indica força e sugere que o movimento deve prosseguir. Por outro lado, se o preço
não consegue alcançar a banda superior, isso pode ser um sinal de fraqueza. Da mesma
forma, a distância do preço em relação às bandas pode revelar informações importantes
sobre a intensidade dos movimentos de alta ou baixa \cite{fernandes2014}.

As Bandas de Bollinger (BB) permitem, adicionalmente, a avaliação da volatilidade do
mercado. Através do indicador BandWidth, derivado das BBs, é possível quantificar a
distância entre as bandas, oferecendo uma medida direta da volatilidade do ativo. A
ocorrência de períodos de baixa volatilidade, caracterizados pela aproximação das
bandas (squeeze), pode indicar um potencial transição para um período de alta
volatilidade (expansão), sinalizando oportunidades de negociação \cite{fernandes2014}.
Em sua pesquisa, Bertini (2024) empregou as Bandas de Bollinger para avaliar a
volatilidade de ativos, demonstrando a utilidade do indicador na seleção de ativos
para estratégias de negociação baseadas na identificação de períodos de variação
acentuada de preços.

\subsection{MACD}
O indicador MACD foi desenvolvido por Gerald Appel na década de 1970; sua proposta é
identificar mudanças no momentum do mercado \cite{appel2005}. Composto por duas
linhas, o MACD (Moving Average Convergence Divergence) representa a diferença entre
duas linhas de médias móveis exponenciais, respectivamente, 12 e 26 períodos; em
seguida, a linha que fornece sinal, média de 9 períodos, o indicador age no cruzamento
de todas as três médias.

Como destacado por \citeonline{elder2002}, o MACD é uma ferramenta interessante na
confirmação de tendências e em identificar pontos específicos de reversão. Geralmente
é utilizado em mercados com tendência mais previsível, mas por derivar da MME, pode
gerar falsos sinais em mercados de lateralidade. Por fim, outra característica
interessante é a capacidade de observar diferenças dos resultados MACD e da
movimentação nos preços, é outra aplicação interessante para a criação de Appel.

\section{Modelos Estatísticos}

Os modelos estatísticos representam uma abordagem fundamental na análise de dados,
especialmente em contextos em que se busca compreender padrões, relações e tendências.
Eles oferecem fórmulas estruturadas para interpretar informações históricas e, em
muitos casos, para projetar comportamentos futuros. Ao aplicar princípios matemáticos
e probabilísticos, esses modelos buscam simplificar a complexidade inerente aos
fenômenos observados, permitindo a extração de insights e a tomada de decisões baseada
em evidências quantitativas.

\subsection{ARIMA}
O modelo Autoregressive Integrated Moving Average (ARIMA) é um modelo estatístico
utilizado para análise de séries temporais. O ARIMA permite modelar dados que exibem
tendências e flutuações ao longo do tempo, sendo amplamente aplicado em áreas
financeiras para prever valores futuros com base em dados passados.

A estrutura do ARIMA é composta por três elementos: o componente Autoregressivo (AR),
que modela a influência dos valores passados da série no valor presente; a Média Móvel
(MA), que suaviza os ruídos e flutuações aleatórias presentes nos dados; e o Integrado
(I), responsável por tornar a série estacionária. Uma série estacionária apresenta
propriedades estatísticas constantes ao longo do tempo, ou seja, sua média e variância
não se alteram significativamente. Cada um desses componentes é definido por um
parâmetro: o parâmetro $p$ representa a ordem da autocorrelação, o parâmetro $q$ define
a ordem da média móvel, e o parâmetro $d$ determina o grau de diferenciação
\cite{xing2024}.

\begin{figure}[htb]
\centering
\caption{Fluxo ARIMA}
\label{fig:fluxo-arima}
\includegraphics[width=0.4\textwidth]{fig05-fluxo-arima.png}
\fonte{Autoria própria.}
\end{figure}

Como é possível observar na figura acima, o processo de aplicação do modelo ARIMA
inicia-se com a coleta do histórico de preços do ativo. Em seguida, verifica-se a
estacionariedade da série; se não estiver estável, aplica-se a diferenciação
repetidamente até que a estabilidade seja alcançada. Com a série pronta, o modelo
ARIMA é ajustado, definindo-se suas ordens $p$, $d$ e $q$. Por fim, o modelo treinado
gera as previsões de preço e os resíduos, que representam os erros da previsão e são
importantes para análises futuras.

A verificação da estacionariedade e o subsequente ajuste dos parâmetros $p$, $d$ e $q$
são metodologicamente embasados em testes e critérios estatísticos. A etapa de
verificar a estacionariedade da série é formalmente conduzida por meio do teste de
Dickey-Fuller Aumentado (ADF). Este teste avalia a hipótese nula de que a série
temporal possui uma raiz unitária, sendo o p-valor resultante um indicador crucial: um
p-valor elevado, como o de 0,998773 encontrado inicialmente por
\citeonline{xing2024} para os dados da NVIDIA, impede a rejeição da hipótese nula,
confirmando a não estacionariedade. Já a aplicação da diferenciação repetidamente até
que a estabilidade seja alcançada é, portanto, guiada pelos resultados sequenciais do
teste ADF. No referido estudo, após a primeira diferenciação, um p-valor próximo de
zero indicou que a estacionariedade foi atingida, definindo assim $d=1$.

Subsequentemente, para o ajuste do modelo ARIMA, um passo crucial é a definição de
suas ordens $p$ (autorregressiva) e $q$ (médias móveis). Dada a multiplicidade de
combinações possíveis para estes parâmetros, como o intervalo de 0 a 4 testado por
\citeonline{xing2024}, torna-se necessário um critério objetivo para a seleção do
modelo mais apropriado. O \textit{Critério de Informação de Akaike} (\textit{AIC}) desempenha
precisamente este papel, sendo uma métrica estatística consolidada para avaliar a
qualidade relativa de diferentes modelos. O AIC permite selecionar a configuração que
oferece o melhor compromisso entre a fidelidade do modelo aos dados históricos
(qualidade do ajuste) e sua complexidade (parcimônia), penalizando modelos com excesso
de parâmetros que poderiam levar a overfitting. Assim, a configuração ARIMA que
resulta no menor valor de AIC é a escolhida, pois representa a estrutura mais eficiente
e generalizável, antes que o modelo treinado seja utilizado para gerar as previsões e
os resíduos \cite{xing2024}.

Uma outra variação do modelo ARIMA é o Seasonal Autoregressive Integrated Moving
Average, ou SARIMA, ideal para dados que exibem padrões sazonais. Ao contrário de
séries aleatórias, dados com sazonalidade apresentam comportamentos previsíveis que se
repetem em intervalos regulares, como vendas que aumentam no Natal ou no verão. O
SARIMA, além dos componentes AR, MA e I, incorpora componentes sazonais para capturar
esses padrões recorrentes, tornando-se uma ferramenta mais dimensionada para modelar
fenômenos com sazonalidade intrínseca \cite{araya2024}.

Apesar de sua utilidade, ambos os modelos apresentam limitações importantes. Sua
eficácia é maior na captura de relações lineares, o que dificulta lidar com padrões não
lineares complexos. Além disso, a escolha dos parâmetros ideais é um desafio, exigindo
análise criteriosa e experimentação. Em cenários complexos e não lineares, modelos de
aprendizado de máquina podem ser mais adequados \cite{araya2024,xing2024}.

\subsection{GARCH}
O modelo Generalized Autoregressive Conditional Heteroskedasticity (GARCH) é uma
ferramenta estatística para a análise de séries temporais financeiras, com foco na
modelagem da volatilidade. A volatilidade, como apontado por \citeonline{araya2024},
reflete o grau de variação dos preços de um ativo financeiro ao longo do tempo e é
fundamental para a gestão de riscos e a precificação de derivativos.

O objetivo principal do modelo GARCH é capturar a dependência temporal da
volatilidade, permitindo que ela varie dinamicamente em resposta às condições do
mercado. O modelo padrão, denominado SGARCH (Standard GARCH), considera que a
volatilidade atual é influenciada tanto pelos choques passados quanto pela própria
volatilidade passada, em um processo de retroalimentação como mostrado na imagem
abaixo \cite{araya2024}.

\begin{figure}[htb]
\centering
\caption{Fluxo GARCH}
\label{fig:fluxo-garch}
\includegraphics[width=0.7\textwidth]{fig06-fluxo-garch.png}
\fonte{Autoria própria.}
\end{figure}

No entanto, o modelo SGARCH assume que a resposta da volatilidade aos choques é
simétrica, ou seja, choques positivos e negativos de igual magnitude têm o mesmo
impacto na volatilidade. Essa premissa nem sempre se verifica nos mercados
financeiros, onde choques negativos (más notícias) frequentemente têm um impacto maior
do que choques positivos (boas notícias). Para lidar com essa assimetria, foram
desenvolvidas variações do modelo GARCH, como o EGARCH (Exponential GARCH) e o TGARCH
(Threshold GARCH) \cite{araya2024}.

O modelo EGARCH, por exemplo, permite que a volatilidade responda de forma diferente a
choques positivos e negativos, capturando o chamado efeito alavancagem. Já o modelo
TGARCH utiliza variáveis binárias (dummy variables) para identificar diferentes regimes
de volatilidade, permitindo que o modelo se adapte a diferentes condições de mercado
\cite{araya2024}.

A escolha entre o modelo SGARCH, EGARCH e TGARCH depende das características específicas
dos dados e dos objetivos da análise. \citeonline{araya2024} destaca essa importância
de selecionar o modelo mais adequado com base em métricas estatísticas.

\subsection{Integração ARIMA-GARCH}

A integração dos modelos ARIMA e GARCH representa uma abordagem robusta para a análise
de séries temporais financeiras, permitindo capturar tanto a tendência condicional
quanto a volatilidade condicional de um ativo. Conforme demonstrado por \citeonline{araya2024},
quando ARIMA e GARCH são combinados, os resíduos gerados pelo modelo ARIMA (que
representam os erros da previsão da média condicional) servem como série de entrada
para o modelo GARCH, que modela a variância condicional desses resíduos.


Essa integração é teoricamente fundamentada no trabalho seminal de \citeonline{bollerslev1986},
que estendeu os modelos de heteroscedasticidade condicional autorregressiva (\textit{ARCH})
inicialmente propostos por \citeonline{engle1982}. O modelo GARCH possibilita capturar
a dinâmica da volatilidade, que é crítica para a gestão de risco em mercados voláteis,
como no caso da ação PETR4, que apresenta elevadas flutuações de preço.


Na prática, a integração ARIMA-GARCH pode ser implementada por meio de três abordagens
estratégicas complementares. A primeira é o \textbf{Filtro de Volatilidade}, que restringe
operações a períodos em que a volatilidade prevista pelo GARCH permanece abaixo de um
limiar definido, reduzindo exposição a mercados instáveis. A segunda é a \textbf{Volatility Targeting}, que ajusta o tamanho da posição inversamente à volatilidade prevista,
diminuindo o risco relativo quando a incerteza aumenta. A terceira é o \textbf{Intervalo de Confiança}, que utiliza a volatilidade do GARCH para estabelecer limites superior
e inferior ao redor das previsões do ARIMA, criando uma zona de confiança para a entrada
e saída de operações.


Para aplicações em negociação de ativos voláteis, o \textbf{Filtro de Volatilidade}
emerge como a abordagem mais prudente, já que permite ao sistema evitar períodos de
heterocedasticidade elevada, alinhando-se ao objetivo de investimento seguro em ativos
com elevada variabilidade de preço. Essa estratégia garante que as operações ocorram
predominantemente em condições de mercado mais estáveis, reduzindo perdas relacionadas a
movimentos abruptos de preço.

\section{Trabalhos Relacionados}

Nesta seção são apresentados estudos e pesquisas relacionados ao tema deste trabalho. O
levantamento realizado foi orientado pela busca de pesquisas científicas e tecnológicas
que têm em seus objetivos o desenvolvimento, implementação e análise de ferramentas e
tecnologias para a previsão de séries temporais financeiras, análise de volatilidade e
desenvolvimento de estratégias de negociação algorítmica, com ênfase na aplicação de
modelos estatísticos e de aprendizado de máquina no mercado de ativos. A ferramenta que
serviu de referência para isso foi o Google Acadêmico e Publish or Perish 8, por meio
do qual se buscou mapear as pesquisas dessa natureza circunscritas nos últimos anos.

O estudo de \citeonline{almeida2016} figura como um precursor direto e ponto de partida
conceitual para o presente trabalho. Almeida investigou a aplicação de métodos
matemáticos e estatísticos no desenvolvimento de um robô de negociação, utilizando a
plataforma MetaTrader 5 e a linguagem MQL5, com foco específico no ativo PETR4. Sua
abordagem centrou-se na estratégia de cruzamento de médias móveis e na análise de
séries temporais com dados históricos, validando o robô através de backtests em
diferentes cenários, incluindo períodos de crise, e com o objetivo de alcançar um
retorno superior à taxa SELIC. O robô demonstrou viabilidade, apresentando
lucratividade em ativos de proteção e superando a SELIC com PETR4 fora de cenários de
crise, além de sugerir a exploração futura de modelos como o GARCH. A pesquisa de
Almeida é de fundamental importância para este TG, pois não apenas validou a utilização
do MetaTrader 5 e MQL5 para a análise algorítmica de PETR4, mas também serviu como base
para a ``revisitação e aplicação contemporânea'' que propomos, conforme mencionado na
justificativa deste projeto. Sua recomendação para explorar modelos mais robustos, como
o GARCH, alinhou-se com a intenção inicial deste trabalho de integrar também redes LSTM.
Embora esse componente tenha sido posteriormente excluído do escopo por limitações de
tempo e de ambiente de desenvolvimento (ver Capítulo 2, seção Obstáculos Encontrados), a
abordagem ARIMA-GARCH efetivamente implementada permitiu construir sobre os achados
iniciais de Almeida, buscando maior capacidade preditiva e adaptabilidade às atuais e
complexas condições do mercado brasileiro.

\citeonline{batista2024} desenvolveu um robô de investimento (\textit{Expert Advisor}) para
automatizar operações na bolsa de valores brasileira (B3), utilizando uma combinação
dos indicadores técnicos Média Móvel Exponencial, Convergência-Divergência de Médias
Móveis e Índice de Força Relativa visando lucros automatizados a médio e longo prazo.
Para isso, empregou a plataforma MetaTrader 5 e a linguagem MQL5, implementando uma
estratégia que integrava os sinais dos três indicadores para decisões de compra e
venda. A validação ocorreu por meio de backtests com dados históricos do ativo BBAS3F
de 2023, comparando o desempenho dos indicadores isoladamente, em pares e em conjunto.
Os resultados mostraram que a combinação dos três indicadores superou as abordagens com
menos indicadores, demonstrando a eficiência e consistência do EA nos testes e
indicando potencial para lucros automatizados.

Em 2022, \citeonline{silva2022} investigaram a previsão de tendências de preços de ações
da Petrobras (PETR4) no mercado brasileiro, com foco em operações intraday de 15
minutos, utilizando técnicas de inteligência artificial. Para isso, empregaram um
extrator de variáveis baseado em \textit{eXtreme Gradient Boosting} (\textit{XGBoost}) para selecionar os
indicadores técnicos mais relevantes dentre mais de 100 disponíveis, visando reduzir
ruído na amostra. Em seguida, uma rede neural \textit{Long Short-Term Memory} (\textit{LSTM}) foi aplicada
para prever a tendência de alta ou baixa dos preços no período de 2014 a 2020. Os
resultados alcançaram uma acurácia de 56\%, superando estudos anteriores, com melhor
desempenho na previsão de quedas (57\%) em comparação com altas (51\%), e destacando a
eficácia da combinação de \textit{XGBoost} para seleção de \textit{features} e \textit{LSTM} para previsão no
contexto do mercado financeiro brasileiro.

Em 2024, \citeonline{xing2024} realizaram uma análise comparativa de quatro modelos de
previsão -- ARIMA, \textit{Multilayer Perceptron} (\textit{MLP}), \textit{Long Short-Term Memory} (\textit{LSTM}) e um
modelo híbrido ARIMA-GARCH -- para prever o preço de fechamento ajustado do dia
seguinte das ações da NVIDIA. Utilizando dados diários da \textit{API} do \textit{Yahoo Finance} de abril
de 2019 a abril de 2024, o estudo dividiu o \textit{dataset} em 90\% para treino e 10\% para
teste. Os modelos \textit{LSTM} e \textit{MLP} foram configurados com janelas deslizantes, enquanto os
modelos ARIMA e ARIMA-GARCH utilizaram janelas expansivas e rolantes, com parâmetros
otimizados para o menor \textit{AIC}. Os resultados, avaliados por \textit{MAE}, \textit{RMSE} e \textit{R}-quadrado,
indicaram que o modelo ARIMA-GARCH apresentou o menor \textit{RMSE}, sugerindo melhor precisão na
previsão de grandes erros e na explicação da variância. No entanto, o modelo ARIMA com
janela expansiva demonstrou desempenho superior em geral, superando os modelos \textit{LSTM} e
\textit{MLP}, possivelmente devido às características predominantemente lineares da série temporal
e à sensibilidade dos modelos de \textit{deep learning} à configuração e ao tamanho da amostra.

\citeonline{araya2024} propuseram um método híbrido GARCH e \textit{Deep Learning} para predição
de volatilidade de preços de ações, visando aumentar a acurácia das previsões ao
integrar diversos modelos. O estudo combinou quatro abordagens: \textit{SARIMA} (\textit{Seasonal
Autoregressive Integrated Moving Average}), modelos da família GARCH (\textit{Generalized
Autoregressive Conditional Heteroskedasticity}), \textit{CNN} (\textit{Convolutional Neural Network}) e
\textit{BiLSTM} (\textit{Bidirectional Long Short-Term Memory}). Inicialmente, os resíduos gerados pelo
modelo \textit{SARIMA} foram utilizados como entrada para o modelo GARCH. Em seguida, a
volatilidade estimada por essa combinação serviu como \textit{feature} de entrada para os
modelos híbridos \textit{CNN} e \textit{BiLSTM}. A \textit{performance} do novo modelo híbrido proposto
(\textit{SARIMA-GARCH-CNN-BiLSTM}) foi avaliada utilizando métricas como \textit{MAE} (\textit{Mean Absolute
Error}) e \textit{RMSE} (\textit{Root Mean Squared Error}), demonstrando uma redução média significativa
nesses erros (60,35\% e 60,61\%, respectivamente) em comparação com outros modelos
híbridos. Os resultados experimentais indicaram que o modelo proposto apresentou bom
desempenho e acurácia na predição da volatilidade de preços de ações, oferecendo
\textit{insights} valiosos para análise de dados financeiros e estratégias de gerenciamento de
risco.

\citeonline{tormin2024} realizou um estudo comparativo entre os modelos ARIMA e \textit{LSTM}
para prever os retornos do Bitcoin, investigando o impacto da integração de dados de
tendências de busca do Google (\textit{PyTrends}) na capacidade preditiva. Utilizando dados
diários de preços do Bitcoin e tendências de busca para o termo ``Bitcoin'' de janeiro
de 2018 a outubro de 2024, foram construídos três modelos: ARIMA, \textit{LSTM} tradicional e
\textit{LSTM} integrado com \textit{PyTrends}. O pré-processamento incluiu ajuste ao fuso horário \textit{UTC},
tratamento de valores ausentes (\textit{fillna}) e normalização (\textit{MinMaxScaler}), com divisão de
80\% dos dados para treino e 20\% para teste. A avaliação, baseada em \textit{RMSE}, \textit{MAE} e \textit{MAPE},
mostrou que o ARIMA, apesar de simples, apresentou maior estabilidade em \textit{RMSE} e \textit{MAE}. O
\textit{LSTM} tradicional obteve \textit{RMSE} levemente melhor que o ARIMA, mas com \textit{MAPE}
significativamente mais alto. A inclusão dos dados do \textit{Google Trends} no modelo \textit{LSTM}
integrado não resultou em melhora significativa, sugerindo a necessidade de ajustes
mais robustos ou a inclusão de outras variáveis para capturar a complexidade do mercado
de criptomoedas.

\begin{table}[htb]
\centering
\caption{Tabela Comparativa}
\label{tab:comparativa}
\fontsize{10}{12}\selectfont
\begin{tabular}{lccc}
\toprule
\textbf{Estudo} & \textbf{Indicadores} & \textbf{Modelo estatístico} & \textbf{Machine Learning} \\
\midrule
(PEREIRA, 2025)         & X & X & X \\
(ALMEIDA, 2016)         & X & X &   \\
(BATISTA, 2024)         & X &   &   \\
(ARAYA, 2024)           &   & X & X \\
(SILVA, 2022)           & X &   & X \\
(XING, 2024)            &   & X & X \\
(TORMIN, 2024)          & X & X & X \\
\bottomrule
\end{tabular}

\vspace{2pt}
{\fontsize{10}{12}\selectfont Comparação entre as pesquisas, evidenciando quais
componentes cada uma utilizou para formular o modelo que elas propuseram.}
\fonte{Autoria própria.}
\end{table}

\begin{table}[htb]
\centering
\caption{Tabela Comparativa (Detalhada)}
\label{tab:comparativa-det}
\fontsize{9}{11}\selectfont
\begin{tabular}{p{2.0cm}cccp{1.6cm}p{1.8cm}p{1.4cm}p{1.8cm}}
\toprule
\textbf{Autor} & \textbf{Ativo Principal} & \textbf{Fonte de Dados} & \textbf{Tempo Gráfico} & \textbf{Indicadores} & \textbf{Modelo(s) Estatístico(s)} & \textbf{Modelo(s) de ML} & \textbf{Métricas para ML} \\
\midrule
(PEREIRA, 2025) & PETR4 & MetaTrader 5 & D1 & RSI, MME, MACD, BB & ARIMA, GARCH & LSTM & MAE, RMSE, R Square \\
(ALMEIDA, 2016) & PETR4 & MetaTrader 5 & D1 & MA & ARIMA, Encontro de Médias Móveis & -- & -- \\
(BATISTA, 2024) & BBAS3F & MetaTrader 5 & M30 & RSI, MME, MACD, BB & -- & -- & -- \\
(SILVA, 2022) & PETR4 & -- & M15 & 113 indicadores & -- & LSTM + XGBoost & Acurácia, Precision, Recall, F1-Score, Support \\
(XING, 2024) & NVIDIA (NVDA) & -- & -- & -- & ARIMA, GARCH & LSTM, MLP & MAE, RMSE, R Square \\
(ARAYA, 2024) & Apple Inc. (AAPL) & Yahoo Finance & D1 & -- & SARIMA, GARCH & CNN, Bi-LSTM & MAE, RMSE \\
(TORMIN, 2024) & Bitcoin (BTC-USD) & Google Trends & -- & -- & ARIMA & LSTM & MAE, RMSE, MAPE \\
\bottomrule
\end{tabular}

\vspace{2pt}
{\fontsize{10}{12}\selectfont Comparação entre as pesquisas, evidenciando as
características de cada uma em diversos aspectos.}
\fonte{Autoria própria.}
\end{table}
